\documentclass[15pt]{scrartcl}
\usepackage[sexy]{evan}
\usepackage{quiver}
\usepackage{relsize}
\newcommand{\PP}{\mathbb{P}}

\begin{document}
\title{Note on \textit{A course in Arithmetic} by J.-P. Serre}
\author{Jiawen Huang}
\maketitle
\newpage
\tableofcontents
\newpage

\section{Chapter 1}
This is a list of proofs of nontrivial results in Chapter 1 of \textit{A course in Arithmetic} by J.-P. Serre.
\\

Let $K$ be a finite field with $q$ elements. 
\begin{theorem}
    Let $u$ be an integer $\geq 0$. The sum 
    \[S(X^u)=\sum_{x\in K} x^u=\begin{cases}
        -1 & \text{if } u>0 \text{ and } (q-1)\mid u\\
        0 & \text{otherwise}
    \end{cases}\]
$u=0$ is subtle here. We define $x^0=1$ for all $x\in K$ (this definition is nice when we work with polynomials because we can write $1$ as $x^0$ to maintain symmetry), so $S(X^0)=q=0$.
\end{theorem}
The proof is easy. Use primitive roots. 
\\
The variant proof: $x\mapsto x^u$ is a group homomorphism from $K^\times $ to itself. The kernel is all $x$ such that $x$ is a root of $x^{\gcd(u,q-1)}-1$. So the image has size $\frac{q-1}{\gcd(u,q-1)}$, which implies the image is just the $d=\frac{q-1}{\gcd(u,q-1)}$ roots of unity. So the sum is $0$ unless $d=1$ (sum of d-th roots of unity is $0$ if $d\geq 2$ s coprime to $p$).
\begin{theorem}[Chevalley-Warning Theorem]
    Let $P_1,\ldots,P_r$ be polynomials in $n$ variables over a finite field $K$ of characteristic $p$. If the sum of their degrees is less than $n$, then the number of common solutions to the equations $P_1=\cdots=P_r=0$ is divisible by $p$. 
\end{theorem}
\begin{proof}
This uses a common trick in counting roots over finite fields. Consider the polynomial
\[Q=\prod_{i=1}^r (1-P_i^{q-1})\]
The number of common solutions is equal to $\sum_{x\in K^n} Q(x)$ since if $P_i(x)=0$ for all $i$, then $Q(x)=1$, otherwise $Q(x)=0$. 
\\

Expand this. For each monomial in the expansion $x^{u_1}_1\cdots x^{u_n}_n$, since $\sum \deg(P_i)(q-1)<n(q-1)$, we have $u_j$ is not divisible less than $q-1$ for some $j$. So by the previous theorem, the sum over $K^n$ of this monomial is $0$. (that's why we define $x^0=1$ above). So the number of common solutions is $0$ mod $p$.
\end{proof}
\begin{corollary}
    All quadratic forms in at least 3 variables over a finite field have nontrivial zeros. 
\end{corollary}
\begin{theorem}
    $\left( \frac{2}{p}\right)=(-1)^{\frac{p^2-1}{8}}$
\end{theorem}
\begin{proof}
    Let $\alpha$ be the primitve 8th root of unity $\alpha$ in $\ol{\FF_p}$. Then, $y^2=(\alpha+\alpha^{-1})^2=2$. If $p\equiv \pm 1\mod 8$, then $y^p=y$, so $2$ is a square in $\FF_p$. If $p\equiv \pm 3\mod 8$, then $y^p=-y$, so $2$ is not a square in $\FF_p$.
\end{proof}
\begin{theorem}[Law of Quadratic Reciprocity]
    For two odd primes $p$ and $q$, we have
    \[\left( \frac{p}{q}\right)\left( \frac{q}{p}\right)=(-1)^{\frac{(p-1)(q-1)}{4}}\]
    
\end{theorem}
The key is to consider the Gauss sum: Let $\zeta$ be a primitive $p$-th root of unity in $\ol{\FF_p}$, then define
\[y=\sum_{x\in \FF_p}\left( \frac{x}{p}\right) \zeta^x\] 
Then we have $y^2=\left( \frac{-1}{p}\right)p=(-1)^{\frac{p-1}{2}}p$ and $y^{q-1}=\left( \frac{q}{p}\right)$. Combining these two gives the result.
\section{Meeting 1 on Jan. 28 2026}
This is a list of nontrivial results in the first meeting of the course.

\begin{theorem}
    The finite subgroup of a field's multiplicative group is cyclic.
\end{theorem}
This can be easily proven using the structure theorem of finite abelian groups and the fundamental theorem of Algebra. 
\\

\begin{definition}
    For any prime $p$ or $p=\infty$, we define the 
    \[\norm{x}_p=\begin{cases}
    p^{-v_p(x)} & p \text{ is a prime}\\
    \norm{x} & p=\infty
    \end{cases}\]
\end{definition}
\begin{theorem}
    For any rational number $r$, $\prod_{p}\norm{r}_p=1$ 
\end{theorem}
This is trivial from the definition of $v_p$.
\\

We can in general define a norm on a number field $K$ as follows: 
\begin{definition}
    $\norm{\cdot}: K\to \RR_{\geq 0}$ is a norm if
    \begin{enumerate}
        \item $\norm{x}=0$ iff $x=0$
        \item $\norm{xy}=\norm{x}\norm{y}$
        \item $\norm{x+y}\leq \norm{x}+\norm{y}$
    \end{enumerate}
\end{definition}
\begin{corollary}
    $\norm{1}=1$ and $\norm{-1}=1$
\end{corollary}
\begin{definition}
    Norm $N_1$ and $N_2$ are equivalent if there exists a real number $c$ such that $N_2(x)=N_1(x)^a$ for all $x$.
\end{definition}
\begin{theorem}[Ostrowski's Theorem]
    Any non-trivial norm on $\QQ$ is equivalent to either the usual absolute value or to $\norm{x}_p$ for some prime $p$.
\end{theorem}
\begin{proof}

\end{proof}

\begin{theorem}[Hensel's Lemma]
\end{theorem}
\begin{proof}
    
\end{proof}
\begin{theorem}[Hasse principle]
    
\end{theorem}
\begin{example}
    $f(x,y,z)=5x^2+7y^2-13z^2$. We want to find a nontrivial $f(x,y,z)=0$ has non-trivial solutions in $\QQ$ using Hensel's Lemma. 
\end{example}
Consider this curve in $\PP^3$, is there any other solution? 
\\

Use parameterization of this curve. 
\\

\begin{theorem}
    If $n\geq 3$, $x^n+y^n=z^n$ has no non-trivial solutions in $\CC[x]$. 
\end{theorem}
\end{document}